%% ---------------------------------------------------------------------------
%% intro.tex
%%
%% Introduction
%%
%% $Id: intro.tex 1477 2010-07-28 21:34:43Z palvarado $
%% ---------------------------------------------------------------------------

\chapter{Introducción}
\label{chp:intro}

\section{Contexto}

El presente proyecto se enmarca en el proyecto HAERA (Hybrid AERial Aquatic robotic system for sampling, monitoring and intervention) de la Universidad de Sevilla, 
que busca avanzar hacia un sistema robótico híbrido aéreo-acuático para muestreo, monitorización e intervención ambiental\cite{Cite2}. Se desarrolla de forma presencial en la 
Escuela Técnica Superior de Ingeniería (ETSI) de dicha universidad \cite{Cite1}, bajo la asesoría de profesores-tutores del Departamento de Ingeniería de Sistemas y Automática, 
en colaboración con el grupo de investigación Multi-Robot and Control Systems (MACS)\cite{Cite3}, especializado en cooperación multi-robot, planificación y control de robots, 
y control no lineal, robusto y adaptativo, con experiencia previa en proyectos multi-robot internacionales y nacionales orientados a aplicaciones medioambientales 
y de inspección.

Para su ejecución, la ETSI dispone de las instalaciones de laboratorio del grupo MACS, que proveen el equipamiento y los recursos necesarios para la integración y 
validación de los sistemas desarrollados, entre ellos las plataformas USV (Unmanned Surface Vehicles) base sobre las que se construye la solución \cite{Cite3}. Las pruebas de 
campo se realizan en cuerpos de agua cercanos a la ciudad de Sevilla, entorno natural para la validación del sistema bajo condiciones reales de operación.

\newpage

\section{Descripción del Problema}

La contaminación de ríos, lagos y océanos, originada por asentamientos humanos, actividades agrícolas e industriales, continúa degradando los ecosistemas acuáticos 
sin que existan, en la mayoría de los casos, herramientas fiables para su monitorización en tiempo real. Esta carencia adquiere especial relevancia considerando que 
más de un tercio de la población mundial carece de acceso a saneamiento mejorado, lo que refuerza la necesidad de datos ambientales de calidad para respaldar la 
gestión de recursos hídricos y la detección temprana de contaminantes\cite{Cite5}.

La ausencia de mecanismos de monitorización sostenida y en tiempo real tiene consecuencias directas sobre la salud de los ecosistemas acuáticos: sin un monitoreo 
adecuado a escala, el daño ambiental resultante puede volverse permanente e irreversible, mientras los cuerpos de agua continúan actuando como sumidero y conducto de 
contaminación antropogénica sin que se generen los datos necesarios para intervenir a tiempo \cite{Cite6}. Esta limitación retrasa la identificación temprana de contaminantes 
y reduce la capacidad de los organismos responsables para tomar decisiones de gestión basadas en información actualizada.

Los sistemas robóticos disponibles actualmente para atender esta necesidad presentan limitaciones documentadas. Los drones aéreos, si bien permiten cubrir grandes 
áreas, alcanzan tasas de éxito de muestreo de entre el 60 {\%} y el 83 {\%} debido a su incapacidad de interactuar directamente con el agua \cite{Cite4}. Por su parte, los vehículos 
aéreo-acuáticos híbridos existentes presentan un tamaño considerable, un costo elevado y una maniobrabilidad reducida, lo que restringe su aplicabilidad en cuerpos de 
agua de difícil acceso. Estas limitaciones evidencian un vacío en la literatura respecto a plataformas multi-robot validadas en condiciones acuáticas no controladas\cite{Cite2}.

\section{Síntesis del Problema}

En el marco del proyecto HAERA, se dispone de plataformas USV base desarrolladas en el laboratorio. Sin embargo, estas carecen de los subsistemas necesarios para operar 
como una flota autónoma orientada a la monitorización y el muestreo de datos en entornos acuáticos reales, limitación documentada como un vacío en plataformas multi-robot 
validadas en condiciones no controladas \cite{Cite2}. Del análisis de las causas asociadas a esta limitación se priorizaron tres causas críticas: la falta de algoritmos de navegación 
autónoma por "waypoints" integrados en la flota, la ausencia de sensores de percepción específicos para misiones de monitorización ambiental acuática, y la falta de 
integración de las plataformas USV base con hardware embarcado modular. Estas tres causas constituyen el eje sobre el cual se estructuran los objetivos del presente proyecto.

El desarrollo del proyecto se enmarca en un periodo de ejecución limitado, condicionado a la adaptación de las plataformas USV ya disponibles en el laboratorio, a la 
variabilidad ambiental de los cuerpos de agua cercanos a Sevilla y al cumplimiento de la normativa española aplicable. Como supuesto crítico, se asume que dichas plataformas 
son adaptables a la integración requerida y que la infraestructura del laboratorio del grupo MACS estará disponible durante toda la ejecución.

Resolver esta limitación permitirá ampliar la cobertura espacial, aumentar la redundancia operativa y mejorar la fiabilidad del muestreo en cuerpos de agua reales. Los 
datasets recolectados en condiciones reales constituirán además un recurso de alto valor para el grupo MACS y para grupos de investigación afines, sentando las bases para 
que organismos de gestión ambiental e instituciones públicas dispongan a futuro de herramientas más confiables para la monitorización autónoma de cuerpos de agua.



\section{Objetivos y estructura del documento}

\index{objetivos}

\subsection{Objetivo General}

Diseñar un sistema multi-USV autónomo que permita ejecutar misiones cooperativas de recolección de datos ambientales en entornos acuáticos.

\subsection{Objetivos Específicos}

\begin{itemize}
    \item Analizar el estado actual de las plataformas USV del laboratorio del grupo MACS para definir los requerimientos oficiales del sistema multi-USV.
    \item Diseñar la arquitectura de integración del hardware embarcado y de los sensores de percepción en al menos dos USVs para habilitar la operación autónoma 
          estable de la plataforma y la adquisición continua de datos ambientales en condiciones controladas de laboratorio.
    \item Implementar la arquitectura de comunicaciones y navegación autónoma cooperativa de la flota multi-USV para habilitar la coordinación simultánea de al menos dos vehículos.
    \item Evaluar el desempeño cooperativo de la flota multi-USV mediante pruebas de integración en un entorno de simulación y en el laboratorio del grupo MACS.
\end{itemize}

\newpage

\subsection{Estructura del documento}

A continuación, se presenta un esquema de las secciones que conforman este documento, junto con una breve descripción de su contenido.

\begin{itemize}
    \item Capítulo 1: Se introduce el contexto del proyecto, definiendo el problema a resolver y los objetivos planteados.
    \item Capítulo 2: Se presenta el marco teórico, con los conceptos fundamentales relevantes para el diseño e implementación del sistema multi-USV.
    \item Capítulo 3: Se expone la metodología de diseño empleada, justificando los procedimientos seguidos para la selección de las alternativas técnicas del sistema.
    \item Capítulo 4: Se detalla la propuesta de diseño, documentando la arquitectura del sistema multi-USV y la integración de sus subsistemas.
    \item Capítulo 5: Se presentan y analizan los resultados obtenidos en las pruebas de validación en simulación y en laboratorio, junto con el análisis económico del proyecto.
    \item Capítulo 6: Se presentan las conclusiones y recomendaciones derivadas del desarrollo del proyecto.

\end{itemize}

El principal aporte de ingeniería de este proyecto radica en la integración y validación de una arquitectura multi-USV autónoma y cooperativa, donde se combinan 
hardware embarcado, sensores de percepción, comunicaciones inalámbricas y navegación por "waypoints" sobre plataformas acuáticas reales, contribuyendo a cerrar el 
vacío identificado en la literatura respecto a sistemas multi-robot validados en condiciones acuáticas no controladas.